\documentclass[../../main.tex]{subfiles}% 注意这里的文件路径不能用 ./main.tex ,否则用latexmk编译子文件会报错
\graphicspath{{\subfix{./image/}}{\subfix{../../image/}}} % 指定图片引用路径,后续可以直接使用图片文件名
% 如果图片存放在上级目录,则使用 ../../image/ 作为备用路径

% 例如:
% \begin{figure}[H]
% \centering
% \includegraphics[width=0.3\textwidth]{图.png}
% \caption{}
% \label{figure:图}
% \end{figure}
% 注意:上述\label{}一定要放在\caption{}之后,否则引用图片序号会只会显示??.

\begin{document}

\section{测度收敛}

\begin{definition}
\begin{enumerate}[(1)]
\item 设 \(1 \leqslant p \leqslant \infty\). 称 \(L^p(X, \mu)\) 中的函数序列 \(f_n\) 强收敛到 \(f\), 如果满足
\begin{align*}
\lim\limits_{n \to \infty} \|f_n - f\|_{L^p} = 0.
\end{align*}
也记为\(f_n \to f\) 在 \(L^p(X, \mu)\)中.

\item 设 \(0 < p < \infty\). 
称序列 \(f_n\) 在 \(L^{p,\infty}(X, \mu)\) 中收敛到 \(f\), 如果满足
\begin{align*}
\lim\limits_{n \to \infty} \|f_n - f\|_{L^{p,\infty}} = 0.
\end{align*}
也记为\(f_n \to f\) 在\(L^{p,\infty}(X, \mu)\) 中.
\end{enumerate}
\end{definition}

\begin{definition}
设 \(f, f_n, n=1, 2, \cdots\) 是测度空间 \((X,\mu)\) 上的可测函数. 如果对于所有 \(\varepsilon > 0\), 存在 \(n_0 \in \mathbb{Z}^+\) 使得
\begin{align*}
n > n_0 \implies \mu(\{x \in X : |f_n(x) - f(x)| > \varepsilon\}) < \varepsilon. 
\end{align*}
则称序列 \(f_n\) \textbf{依测度收敛}到 \(f\). 这就等价于
\begin{align*}
\lim\limits_{n \to \infty} \mu(\{x \in X : |f_n(x) - f(x)| > \varepsilon\}) = 0. 
\end{align*}
\end{definition}
\begin{remark}
依测度收敛是一个比 \(L^p\) 或 \(L^{p,\infty}\) 收敛更弱的概念.
\end{remark}

\begin{proposition}
设 \(0 < p \leqslant \infty\) 且 \(f_n, f\) 属于 \(L^p(X,\mu)\).
\begin{enumerate}[(1)]
\item 若 \(f_n, f \in L^p\) 且 \(f_n \to f\) 在 \(L^p\) 中, 则 \(f_n \to f\) 在 \(L^{p,\infty}\) 中.
\item 若 \(f_n \to f\) 在 \(L^{p,\infty}\) 中, 则 \(f_n\) 依测度收敛到 \(f\).
\end{enumerate}
\end{proposition}
\begin{proof}
\begin{enumerate}[(1)]
\item  固定 \(0 < p < \infty.\) 对于所有 \(\alpha > 0,\) 由\hyperref[Real Analysis-theorem:实分析-Chebyshev不等式]{Chebyshev不等式},我们有
\begin{align*}
\mu(\{x \in X : |f_n(x) - f(x)| > \alpha\}) \leqslant \frac{1}{\alpha^p} \int_X |f_n - f|^p \mathrm{d}\mu.
\end{align*}
从而
\begin{align*}
\parallel f_n-f\parallel _{L^{p,\infty}}=\mathop {\mathrm{sup}} \limits_{\alpha >0}\alpha \mu (\{x\in X:|f_n(x)-f(x)|>\alpha \})^{\frac{1}{p}}<\mathop {\mathrm{sup}} \limits_{\alpha >0}\left\{ \alpha \cdot \frac{\left\| f_n-f \right\| _{L^p}}{\alpha} \right\} =\left\| f_n-f \right\| _{L^p}.
\end{align*}
这表明 \(L^p\) 中的收敛蕴含着弱 \(L^p\) 中的收敛. 当 \(p = \infty\) 时,由定义知$L^{\infty,\infty}=L^{\infty}$,故此时结论显然成立.

\item 给定 \(\varepsilon > 0,\) 存在 \(n_0\) 使得对于 \(n > n_0\) 有
\begin{align*}
\|f_n - f\|_{L^{p,\infty}} = \sup\limits_{\alpha > 0} \alpha \mu(\{x \in X : |f_n(x) - f(x)| > \alpha\})^{\frac{1}{p}} < \varepsilon^{\frac{1}{p}+1}.
\end{align*}
取 \(\alpha = \varepsilon,\) 可得 \(L^{p,\infty}\) 中的收敛蕴含着依测度收敛.
\end{enumerate}
\end{proof}

\begin{example}
固定 \(0 < p < \infty.\) 在 \([0, 1]\) 上定义函数
\begin{align*}
f_{k,j} = k^{\frac{1}{p}} \chi_{\left(\frac{j-1}{k}, \frac{j}{k}\right)}, \quad k \geqslant 1, 1 \leqslant j \leqslant k.
\end{align*}
考虑序列 \(\{f_{1,1}, f_{2,1}, f_{2,2}, f_{3,1}, f_{3,2}, f_{3,3}, \cdots\}.\) 观察得到
\begin{align*}
|\{x : f_{k,j}(x) > 0\}| = \frac{1}{k}.
\end{align*}
因此, \(f_{k,j}\) 依测度收敛到 \(0.\) 同样地, 观察可得
\begin{align*}
\|f_{k,j}\|_{L^{p,\infty}} = \sup\limits_{\alpha > 0} \alpha |\{x : f_{k,j}(x) > \alpha\}|^{\frac{1}{p}} \geqslant \sup\limits_{k \geqslant 1} \frac{\left(\frac{k - 1}{k}\right)^{\frac{1}{p}}}{k^{\frac{1}{p}}} = 1,
\end{align*}
这意味着 \(f_{k,j}\) 在 \(L^{p,\infty}\) 中不收敛到 \(0.\)
\end{example}

\begin{theorem}
设 \(f_n\) 和 \(f\) 是测度空间 \((X, \mu)\) 上的复值可测函数, 且假设 \(f_n\) 依测度收敛到 \(f.\) 则存在 \(f_n\) 的一个子列使得 \(f_{n_k}\) 依测度 \(\mu\)-a.e. 收敛到 \(f.\)
\end{theorem}

\begin{proof}
对于所有的 \(k = 1, 2, \cdots,\) 归纳地选取 \(n_k\) 使得
\begin{align}
\mu(\{x \in X : |f_{n_k}(x) - f(x)| > 2^{-k}\}) < 2^{-k}, \label{eq:1.1.17}
\end{align}
且满足 \(n_1 < n_2 < \cdots < n_k < \cdots.\) 定义集合
\begin{align*}
A_k = \{x \in X : |f_{n_k}(x) - f(x)| > 2^{-k}\}.
\end{align*}
由 \eqref{eq:1.1.17} 可知
\begin{align}
\mu\left( \bigcup\limits_{k=m}^{\infty} A_k \right) \leqslant \sum\limits_{k=m}^{\infty} \mu(A_k) \leqslant \sum\limits_{k=m}^{\infty} 2^{-k} = 2^{1-m}, \label{eq:1.1.18}
\end{align}
对于所有 \(m = 1, 2, 3, \cdots\) 成立. 由 \eqref{eq:1.1.18} 可得
\begin{align}
\mu\left( \bigcup\limits_{k=1}^{\infty} A_k \right) \leqslant 1 < \infty. \label{eq:1.1.19}
\end{align}
利用 \eqref{eq:1.1.18} 和 \eqref{eq:1.1.19}, 我们可得集合序列 \(\{\bigcup\limits_{k=m}^{\infty} A_k\}_{m=1}^{\infty}\) 的测度当 \(m \to \infty\) 时收敛到 \(0,\) 即
\begin{align}
\mu\left( \bigcap\limits_{m=1}^{\infty} \bigcup\limits_{k=m}^{\infty} A_k \right) = 0. \label{eq:1.1.20}
\end{align}
为了完成证明, 注意到 \eqref{eq:1.1.20} 中的零测集包含了使得 \(f_{n_k}(x)\) 不收敛到 \(f(x)\) 的所有 \(x \in X\) 的集合.
\end{proof}

\begin{definition}
我们说测度空间 \((X, \mu)\) 上的可测函数序列 \(\{f_n\}\) 是依测度 Cauchy 列的, 如果对于所有 \(\varepsilon > 0,\) 存在 \(n_0 \in \mathbb{Z}^+\) 使得对于 \(n, m > n_0\) 都有
\begin{align*}
\mu(\{x \in X : |f_m(x) - f_n(x)| > \varepsilon\}) < \varepsilon.
\end{align*}
\end{definition}

\begin{theorem}
设 \((X, \mu)\) 是测度空间, 且 \(\{f_n\}\) 是 \(X\) 上的依测度 Cauchy 列. 则存在 \(f_n\) 的一个子列在 \(X\) 上 \(\mu\)-a.e. 收敛.
\end{theorem}

\begin{proof}
该证明与定理 1.1.11 的证明非常类似. 对于所有 \(k = 1, 2, \cdots,\) 归纳地选取 \(n_k\) 使得
\begin{align}
\mu(\{x \in X : |f_{n_k}(x) - f_{n_{k+1}}(x)| > 2^{-k}\}) < 2^{-k}, \label{eq:1.1.21}
\end{align}
且满足 \(n_1 < n_2 < \cdots < n_k < n_{k+1} < \cdots.\) 定义
\begin{align*}
A_k = \{x \in X : |f_{n_k}(x) - f_{n_{k+1}}(x)| > 2^{-k}\}.
\end{align*}
如定理 1.1.11 的证明所示, \eqref{eq:1.1.21} 蕴含
\begin{align}
\mu\left( \bigcap\limits_{m=1}^{\infty} \bigcup\limits_{k=m}^{\infty} A_k \right) = 0. \label{eq:1.1.22}
\end{align}
对于 \(x \notin \bigcup\limits_{k=m}^{\infty} A_k\) 且 \(i \geqslant j \geqslant j_0 \geqslant m\) (其中 \(j_0\) 足够大), 我们有
\begin{align*}
|f_{n_i}(x) - f_{n_j}(x)| \leqslant \sum\limits_{l=j}^{i-1} |f_{n_l}(x) - f_{n_{l+1}}(x)| \leqslant \sum\limits_{l=j}^{i-1} 2^{-l} \leqslant 2^{1-j} \leqslant 2^{1-j_0}.
\end{align*}
这意味着序列 \(\{f_{n_i}(x)\}_i\) 是集合 \((\bigcup\limits_{k=m}^{\infty} A_k)^c\) 中每个 \(x\) 的 Cauchy 列, 因此对于所有这样的 \(x\) 都收敛. 我们定义一个函数
\begin{align*}
f(x) = 
\begin{cases} 
\lim\limits_{j \to \infty} f_{n_j}(x), & \text{当 } x \notin \bigcap\limits_{m=1}^{\infty} \bigcup\limits_{k=m}^{\infty} A_k, \\
0, & \text{当 } x \in \bigcap\limits_{m=1}^{\infty} \bigcup\limits_{k=m}^{\infty} A_k.
\end{cases}
\end{align*}
那么 \(f_{n_j} \to f\) 几乎处处成立.
\end{proof}




















\end{document}